% GENERATED FILE. Edit canonical lessons, then run npm run generate:books.
% canonical-source-hash: fnv1a64:10fdd8ed1d2621aa
% canonical-lessons: MR-C124-yojana, MR-C124-pahat, MR-C124-madhyaratra, MR-C124-parva, MR-C124-somvar

\chapter{Plan, Dawn, Midnight, Day after tomorrow, Monday}
\label{ch:mr-plan-dawn-midnight-day-after-tomorrow-monday}
\hlchaptermodality{124}

\begin{chapteropening}
\textbf{By the end of this chapter:} \emph{I can say a plan, dawn, midnight, the day after tomorrow and Monday.}

Complete the last lesson of chapter 124: I can say a plan, dawn, midnight, the day after tomorrow and Monday.
\end{chapteropening}

\section[\texorpdfstring{\mr{योजना} (yojanā)}{योजना (yojanā)}]{\mr{योजना} (yojanā) — a plan}

\label{lesson:MR-C124-yojana}



\begin{quote}
Before the new one: say the Marathi for to ask for, then the Marathi for to decide.
\end{quote}

\subsection*{You'll want to know: \mr{योजना}}
\textbf{\mr{योजना}} — \emph{yojanā} — \textquotedblleft{}a plan\textquotedblright{}.

\begin{cousinweb}
\textbf{\mr{योजना}} (\emph{yojanā}) is Sanskrit \emph{yojanā}, \textquotedblleft{}joining, arrangement\textquotedblright{}, a tatsama built on \emph{yuj}, \textquotedblleft{}to yoke, join\textquotedblright{}. CDIAL has no entry for it.
\end{cousinweb}

\begin{grammarlens}[title={What you've built}]
One more word for this chapter.
\end{grammarlens}

\subsection*{Your turn}
\begin{itemize}
\raggedright
  \item \emph{Say it:} \emph{yojanā}
  \item \emph{Say it:} \emph{yojanā}, once more
  \item \emph{Say it:} say \emph{yojanā}
  \item \emph{Recall:} say the Marathi for to ask for, then the Marathi for to decide, then say \emph{yojanā} again
\end{itemize}

\begin{tcolorbox}[breakable,colback=teal!4,colframe=teal!35!black,title={Before you move on}]
What does \mr{योजना} mean? (\textquotedblleft{}A plan\textquotedblright{}.) Say it once more.
\end{tcolorbox}
\section[\texorpdfstring{\mr{पहाट} (pahāṭ)}{पहाट (pahāṭ)}]{\mr{पहाट} (pahāṭ) — dawn}

\label{lesson:MR-C124-pahat}



\begin{quote}
Before the new one: say the Marathi for to decide, then the Marathi for a plan.
\end{quote}

\subsection*{You'll want to know: \mr{पहाट}}
\textbf{\mr{पहाट}} — \emph{pahāṭ} — \textquotedblleft{}dawn\textquotedblright{}.

\begin{cousinweb}
\textbf{\mr{पहाट}} (\emph{pahāṭ}) is inherited. CDIAL lists Marathi \emph{pahāṭ}, \textquotedblleft{}period before sunrise, dawn\textquotedblright{}, under Sanskrit \emph{prabhā}, \textquotedblleft{}light\textquotedblright{}, with a query on how the word was formed (CDIAL 8705).
\end{cousinweb}

\begin{grammarlens}[title={What you've built}]
One more word for this chapter.
\end{grammarlens}

\subsection*{Your turn}
\begin{itemize}
\raggedright
  \item \emph{Say it:} \emph{pahāṭ}
  \item \emph{Say it:} \emph{pahāṭ}, once more
  \item \emph{Say it:} say \emph{pahāṭ}
  \item \emph{Recall:} say the Marathi for to decide, then the Marathi for a plan, then say \emph{pahāṭ} again
\end{itemize}

\begin{tcolorbox}[breakable,colback=teal!4,colframe=teal!35!black,title={Before you move on}]
What does \mr{पहाट} mean? (\textquotedblleft{}Dawn\textquotedblright{}.) Say it once more.
\end{tcolorbox}
\section[\texorpdfstring{\mr{मध्यरात्र} (madhyarātra)}{मध्यरात्र (madhyarātra)}]{\mr{मध्यरात्र} (madhyarātra) — midnight}

\label{lesson:MR-C124-madhyaratra}



\begin{quote}
Before the new one: say the Marathi for a plan, then the Marathi for dawn.

Walk through the day before its end: \textbf{\mr{दिवस}}, then the evening greeting, \textbf{\mr{शुभ} \mr{संध्याकाळ}}.
\end{quote}

\subsection*{You'll want to know: \mr{मध्यरात्र}}
\textbf{\mr{मध्यरात्र}} — \emph{madhyarātra} — \textquotedblleft{}midnight\textquotedblright{}.

\begin{cousinweb}
\textbf{\mr{मध्यरात्र}} (\emph{madhyarātra}) is a Sanskrit compound, a tatsama: \emph{madhya}, \textquotedblleft{}middle\textquotedblright{} (CDIAL 9804), and \emph{rātra}, \textquotedblleft{}night\textquotedblright{}.
\end{cousinweb}

\begin{grammarlens}[title={What you've built}]
One more word for this chapter.
\end{grammarlens}

\subsection*{Your turn}
\begin{itemize}
\raggedright
  \item \emph{Say it:} \emph{madhyarātra}
  \item \emph{Say it:} \emph{madhyarātra}, once more
  \item \emph{Say it:} say \emph{madhyarātra}
  \item \emph{Recall:} say the Marathi for a plan, then the Marathi for dawn, then say \emph{madhyarātra} again
\end{itemize}

\begin{tcolorbox}[breakable,colback=teal!4,colframe=teal!35!black,title={Before you move on}]
What does \mr{मध्यरात्र} mean? (\textquotedblleft{}Midnight\textquotedblright{}.) Say it once more.
\end{tcolorbox}
\section[\texorpdfstring{\mr{परवा} (parvā)}{परवा (parvā)}]{\mr{परवा} (parvā) — the day after tomorrow}

\label{lesson:MR-C124-parva}



\begin{quote}
Before the new one: say the Marathi for dawn, then the Marathi for midnight.
\end{quote}

\subsection*{You'll want to know: \mr{परवा}}
\textbf{\mr{परवा}} — \emph{parvā} — \textquotedblleft{}the day after tomorrow\textquotedblright{}.

\begin{cousinweb}
\textbf{\mr{परवा}} (\emph{parvā}) is a tadbhava, a word inherited from Sanskrit and reshaped by sound change. It comes from \emph{paraśvas}, \textquotedblleft{}the day after tomorrow\textquotedblright{} (CDIAL 7799j). It is \emph{para}, \textquotedblleft{}beyond\textquotedblright{}, and \emph{śvas}, \textquotedblleft{}tomorrow\textquotedblright{}. CDIAL glosses the Marathi word \textquotedblleft{}the day after tomorrow or before yesterday\textquotedblright{}.
\end{cousinweb}

\begin{grammarlens}[title={What you've built}]
One more word for this chapter.
\end{grammarlens}

\subsection*{Your turn}
\begin{itemize}
\raggedright
  \item \emph{Say it:} \emph{parvā}
  \item \emph{Say it:} \emph{parvā}, once more
  \item \emph{Say it:} say \emph{parvā}
  \item \emph{Recall:} say the Marathi for dawn, then the Marathi for midnight, then say \emph{parvā} again
\end{itemize}

\begin{tcolorbox}[breakable,colback=teal!4,colframe=teal!35!black,title={Before you move on}]
What does \mr{परवा} mean? (\textquotedblleft{}The day after tomorrow\textquotedblright{}.) Say it once more.
\end{tcolorbox}
\section[\texorpdfstring{\mr{सोमवार} (somvār)}{सोमवार (somvār)}]{\mr{सोमवार} (somvār) — Monday}

\label{lesson:MR-C124-somvar}



\begin{quote}
Before the new one: say the Marathi for midnight, then the Marathi for the day after tomorrow.
\end{quote}

\subsection*{You'll want to know: \mr{सोमवार}}
\textbf{\mr{सोमवार}} — \emph{somvār} — \textquotedblleft{}Monday\textquotedblright{}.

\begin{cousinweb}
\textbf{\mr{सोमवार}} (\emph{somvār}) is a Sanskrit compound, a tatsama: \emph{soma}, \textquotedblleft{}the moon\textquotedblright{}, and \emph{vāra}, \textquotedblleft{}day of the week\textquotedblright{} (CDIAL 11547). Monday is the moon's day, as in English.
\end{cousinweb}

\begin{grammarlens}[title={What you've built}]
One more word for this chapter.
\end{grammarlens}

\subsection*{Your turn}
\begin{itemize}
\raggedright
  \item \emph{Say it:} \emph{somvār}
  \item \emph{Say it:} \emph{somvār}, once more
  \item \emph{Say it:} say \emph{somvār}
  \item \emph{Recall:} say the Marathi for midnight, then the Marathi for the day after tomorrow, then say \emph{somvār} again
\end{itemize}

\begin{tcolorbox}[breakable,colback=teal!4,colframe=teal!35!black,title={Before you move on}]
What does \mr{सोमवार} mean? (\textquotedblleft{}Monday\textquotedblright{}.) Say it once more.
\end{tcolorbox}
